\par\Needspace{6\baselineskip}\medskip
\SourceLeaf{136}\pdfbookmark[2]{98.}{article-098}
\begin{center}98.\end{center}
Solutio problematis nostri pro ellipsi in praecc. evoluta etiam ad parabolam et hyperbolam transferri posset, considerando parabolam tamquam ellipsin, in qua $a$ et $b$ essent quantitates infinitae, $\varphi=90^0$, tandem $E$, $E'$, $g$, $G=0$; et perinde hyperbolam tamquam ellipsin in qua $a$ esset negativa, atque $b$, $E$, $E'$, $g$, $G$, $\varphi$ imaginariae: malumus tamen his suppositionibus\raisebox{.45ex}{.} abstinere, problemaque pro utroque sectionum conicarum genere seorsim tractare. Analogia insignis inter omnia tria genera sic sponte se manifestabit.

Retinendo in PARABOLA characteres $p$, $v$, $v'$, $F$, $f$, $r$, $r'$, $t$ in eadem significatione in qua supra accepti sunt, habemus e theoria motus parabolici:
\[
\begin{aligned}
\mathrm{1)}\qquad&\BareRoot{\dfrac p{2r}}=\cos\tfrac12(F-f)\\[.7em]
\mathrm{2)}\qquad&\BareRoot{\dfrac p{2r'}}=\cos\tfrac12(F+f)
\end{aligned}
\]
\[
\begin{aligned}
\frac{2kt}{p^{\frac32}}&=\operatorname{tang}\tfrac12(F+f)-\operatorname{tang}\tfrac12(F-f)\\
&\quad+\tfrac13\operatorname{tang}\tfrac12(F+f)^3-\tfrac13\operatorname{tang}\tfrac12(F-f)^3\\[.5em]
&=\{\operatorname{tang}\tfrac12(F+f)-\operatorname{tang}\tfrac12(F-f)\}\\
&\quad\Bigl\{1+\operatorname{tang}\tfrac12(F+f)\operatorname{tang}\tfrac12(F-f)\\
&\qquad+\tfrac13(\operatorname{tang}\tfrac12(F+f)-\operatorname{tang}\tfrac12(F-f))^2\Bigr\}\\[.5em]
&=\frac{2\sin f\,\mathpunct{.}\BareRoot{rr'}}p
\left\{\frac{2\cos f\,\mathpunct{.}\BareRoot{rr'}}p+\frac{4\sin f^2rr'}{3pp}\right\},
\end{aligned}
\]
\noindent unde
\[
\mathrm{3)}\qquad kt=\frac{2\sin f\cos f\,\mathpunct{.}\,rr'}{\BareRoot{p}}+\frac{4\sin f^3(rr')^{\frac32}}{3p^{\frac32}}\mathpunct{\cdot}
\]
\noindent Porro deducitur ex multiplicatione aequationum 1, 2
\[
\mathrm{4)}\qquad \frac p{\BareRoot{rr'}}=\cos F+\cos f,
\]
\noindent nec non ex additione quadratorum
\[
\mathrm{5)}\qquad \frac{p(r+r')}{2rr'}=1+\cos F\cos f.
\]
\noindent Hinc eliminato $\cos F$
\[
\mathrm{6)}\qquad p=\frac{2rr'\sin f^2}{r+r'-2\cos f\,\mathpunct{.}\BareRoot{rr'}}\mathpunct{\cdot}
\]
\noindent Quodsi itaque aequationes 9, $9^*$ art. 88 hic quoque adoptamus, priorem pro \SourceLeaf{137}$\cos f$ positivo, posteriorem pro negativo, habebimus
\[
\begin{aligned}
\mathrm{7)}\qquad&p=\frac{\sin f^2\BareRoot{rr'}}{2l\cos f}\\[.6em]
\mathrm{7^*)}\qquad&p=\frac{\sin f^2\BareRoot{rr'}}{-2L\cos f}\mathpunct{\raisebox{.45ex}{,}}
\end{aligned}
\]
\noindent quibus valoribus in aequatione 3 substitutis, prodibit, retinendo characteres $m$, $M$ in significatione per aequationes 11, $11^*$ art. 88 stabilita,
\[
\begin{aligned}
\mathrm{8)}\qquad&m=l^{\frac12}+\tfrac43l^{\frac32}\\[.5em]
\mathrm{8^*)}\qquad&M=-L^{\frac12}+\tfrac43L^{\frac32}.
\end{aligned}
\]
\noindent Hae aequationes conveniunt cum 12, $12^*$ art. 88, si illic statuatur $g=0$. Hinc colligitur, si duo loci heliocentrici, quibus per parabolam satisfit, ita tractentur, ac si orbita esset elliptica, ex applicatione praeceptorum art. 91 statim resultare debere $x=0$; vice versa facile perspicitur, si per praecepta ista prodeat $x=\mathpunct{\raisebox{.65ex}{.}}0$, orbitam pro ellipsi parabolam evadere, quum per aequationes 1, 16, 17, 19, 20 fiat $b=\infty$, $a=\infty$, $\varphi=90^0$. Determinatio elementorum facillime dein absolvitur. Pro $p$ enim adhiberi poterit vel aequatio 7 art. praesentis, vel aequ. 18 art. 95\footnote{Unde simul patet, $y$ et $Y$ in parabola easdem rationes exprimere ut in ellipsi, v. art. 95.}: pro $F$ autem fit ex aequationibus 1, 2 huius art. $\operatorname{tang}\tfrac12 F=\dfrac{\BareRoot{r'}-\BareRoot{r}}{\BareRoot{r'}+\BareRoot{r}}\operatorname{cotang}\tfrac12 f=\sin2\omega\operatorname{cotang}\tfrac12 f$, si angulus auxiliaris in eadem significatione accipitur, ut in art. 89.

Hacce occasione adhuc observamus, si in aequ. 3 pro $p$ substituatur valor eius ex 6, prodire aequationem satis notam
\[
kt=\tfrac13(r+r'+\cos f\,\mathpunct{.}\BareRoot{rr'})(r+r'-2\cos f\,\mathpunct{.}\BareRoot{rr'})^{\frac12}\BareRoot{2}.
\]
